\documentclass[11pt]{article}

\usepackage[margin=1.1in]{geometry}
\usepackage{amsmath,amssymb,amsthm}
\usepackage{booktabs}
\usepackage{listings}
\usepackage{xcolor}
\usepackage{microtype}
\usepackage[hidelinks]{hyperref}
\emergencystretch=1.5em

\newtheorem{theorem}{Theorem}[section]
\newtheorem{proposition}[theorem]{Proposition}
\newtheorem{corollary}[theorem]{Corollary}
\theoremstyle{remark}
\newtheorem{remark}[theorem]{Remark}

\lstset{
  basicstyle=\ttfamily\small,
  columns=fullflexible,
  keepspaces=true,
  breaklines=true,
  literate={∃}{{$\exists$}}1
           {∧}{{$\land$}}1
           {≤}{{$\le$}}1
           {ℝ}{{$\mathbb{R}$}}1
           {ℕ}{{$\mathbb{N}$}}1
           {σ}{{$\sigma$}}1,
}

\newcommand{\tr}{\operatorname{tr}}
\newcommand{\loopsupp}{P}
\newcommand{\Area}{\mathrm{A}}

\title{A Machine-Checked Volume-Uniform Wilson-Loop Area Law\\
via a Formalized Cluster Expansion}
\author{Lluis Eriksson}
\date{July 2, 2026}

\begin{document}
\maketitle

\begin{abstract}
We report a complete formalization, in Lean~4 over Mathlib, of
volume-uniform Wilson-loop area laws for $SU(N_c)$ lattice gauge theory in
an explicit strong-coupling window --- including the case of the
\emph{exact} Wilson Boltzmann factor $\prod_p e^{z_p}$, not a linearized
surrogate.  The headline theorem bounds the normalized Wilson-loop
expectation by
$N_c\, e^{\loopsupp \cdot 4d \cdot K}\,\sigma^{\Area(C)}\,
e^{\loopsupp \cdot 4d \cdot S(\sigma)}$,
where $\Area(C)$ is an intrinsic combinatorial filling area of the loop,
$\loopsupp$ is its edge-support size, and \emph{every constant is
volume-free}: the bound holds uniformly over all finite lattice sizes.  The
partition function is cancelled through a fully formalized volume-restricted
cluster expansion (loop-tagged factorization, restricted Mayer inversion,
$Z$-ratio bounds, and a pinned-gas resummation built on a
Koteck\'y--Preiss layer with Penrose-style spanning-tree counting).  A
reusable repackaging converts the bound into manifest exponential area decay
with a strictly positive string tension, and --- reflecting a methodological
lesson from the project's own history --- the \emph{non-vacuity} of every
hypothesis window is itself machine-checked --- both the cluster smallness
window and the decay-repackaging window, the latter with an explicit
witness of tension $\log 2 - \tfrac12$.  For every exported theorem in this
chain the Lean kernel's axiom oracle reports exactly
\texttt{[propext, Classical.choice, Quot.sound]}; there is no \texttt{sorry}
and no project axiom in the dependency cone.  To our knowledge this is the
first machine-checked cluster-expansion proof in lattice quantum field
theory.  All artifacts are public, with per-theorem oracle records in a
verification ledger and continuous-integration builds.
\end{abstract}

\section{Introduction}\label{sec:intro}

Wilson's confinement criterion asks whether the expectation of a Wilson
loop decays with the \emph{area} it spans rather than with its
perimeter~\cite{Wilson1974}.  At strong coupling this is a classical
theorem of lattice gauge theory: the character/cluster expansions of
Osterwalder and Seiler establish area-law decay for the Wilson action in a
small-$\beta$ regime~\cite{OsterwalderSeiler1978,Seiler1982}.  The
underlying machinery --- polymer representations, Mayer expansions, and
tree-graph bounds in the tradition of Penrose~\cite{Penrose1967} and
Koteck\'y--Preiss~\cite{KoteckyPreiss1986} --- is standard, but it is also
notoriously bookkeeping-heavy: the mathematical content of
``volume-uniform'' is precisely that \emph{every} constant produced by the
expansion can be tracked through the thermodynamic limit without silent
volume dependence.  This is the kind of proof in which human errors hide in
constants rather than in ideas, and hence a natural target for machine
verification.

This paper reports such a verification.  Within the author's larger
audit-first programme on lattice Yang--Mills
theory~\cite{Eriksson2026Master,Eriksson2026Audit}, the following family of
results has been formalized end-to-end in Lean~4~\cite{Lean4} over the
Mathlib library~\cite{mathlib}:

\begin{enumerate}
\item \textbf{Area laws in four variants} --- finite-volume and
  volume-uniform, each for linearized activities and for the \emph{exact}
  Wilson Boltzmann factor $\prod_p e^{z_p}$ --- with explicit constants and
  explicit smallness windows (Section~\ref{sec:results}).
\item \textbf{A formalized cluster expansion} strong enough to cancel the
  partition function with volume-independent losses: loop-tagged
  factorization, restricted Mayer inversion, $Z$-ratio bounds, and a
  pinned-gas resummation over a Koteck\'y--Preiss layer with Penrose-style
  breadth-first-search spanning-tree counting
  (Section~\ref{sec:architecture}).
\item \textbf{A methodological discipline} in which integrability side
  conditions are theorems rather than hypotheses, per-theorem axiom oracles
  are recorded in a public ledger, and --- after an early episode in which
  the repository itself had advertised a vacuous headline statement, since
  disclaimed in its README --- the non-vacuity of every hypothesis window
  is itself machine-checked (Section~\ref{sec:lean}).
\end{enumerate}

The mathematics is classical in spirit; the contribution is the machine
checking, the explicitness, and the reusable formal toolbox.  We are not
aware of any prior machine-checked cluster or polymer expansion for a
lattice quantum field theory in any proof assistant.

\section{Setting}\label{sec:setting}

Fix a dimension $d \ge 1$ and a linear lattice size $N \ge 1$ (both
positive; the theorems are stated for all such $d, N$).  The repository's
lattice layer provides a concrete complex of oriented edges and plaquettes
over the finite hypercubic lattice of linear size $N$; a gauge
configuration is a map $A$ from oriented edges to the gauge group
$G = SU(N_c)$ (as Mathlib's \texttt{Matrix.specialUnitaryGroup}) satisfying
$A(e^{-1}) = A(e)^{-1}$.  The reference measure $\mu_N$ is the product of
normalized Haar measures over one orientation class of edges, transported
to the configuration space (\texttt{gaugeMeasureFrom},
\texttt{sunHaarProb}).

For a list of edges $\mathit{es}$ tracing a closed loop $C$, the Wilson
line $W_C(A)$ is the ordered product of the edge variables and
$\tr W_C$ its (matrix) trace.  Plaquette activities take the
conjugate-pair form
\[
  z_p(A) \;=\; c_p \cdot \tr H_p(A) \;+\; \overline{c_p} \cdot
  \overline{\tr H_p(A)},
  \qquad \lVert c_p \rVert \le \delta ,
\]
where $H_p$ is the plaquette holonomy; the conjugate pairing makes each
$z_p$ real, and the weight $\prod_p e^{z_p}$ is the \emph{exact} Wilson
Boltzmann factor.  In particular the standard Wilson action weight
$\exp\bigl((\beta/N_c)\operatorname{Re}\tr H_p\bigr)$ is the case
$c_p = \beta/(2N_c)$, so $\delta = \beta/(2N_c)$.  The partition function
is $Z_N = \int \prod_p e^{z_p}\, d\mu_N$, and the object of interest is the
normalized expectation
$\langle \tr W_C \rangle_N = Z_N^{-1} \int \tr W_C \cdot \prod_p e^{z_p}\,
d\mu_N$.

Two combinatorial quantities of the loop control the bound.  The
\emph{perimeter charge} $\loopsupp = \loopsupp(C)$ is the number of
distinct edges supporting the loop (\texttt{edgeSupport}).  The
\emph{area} $\Area(C)$ is intrinsic and algebraic: the loop determines a
$1$-chain with coefficients in $\mathbb{Z}/N_c$, and
\[
  \Area(C) \;=\; \min \bigl\{\, \#\mathrm{supp}(s) \;:\;
    s \text{ a plaquette 2-chain over } \mathbb{Z}/N_c,\;
    \partial_2 s = [C] \,\bigr\}
\]
(\texttt{chainAreaA}, \texttt{loopChain}) --- the minimal number of
plaquettes in an algebraic filling of the loop.  No area hypothesis is
imposed anywhere; the exponent in the theorems below is this defined
quantity.

\section{Main results}\label{sec:results}

Write $\rho(\delta) = e^{2\delta N_c} - 1$ for the effective exact-activity
size, and define the two window functions
\[
  K(t;\delta) \;=\;
    \frac{e \cdot \rho(\delta)\, e^{t}}
         {1 - (16d+1)^2\, \rho(\delta)\, e^{t}},
  \qquad
  S(\sigma) \;=\; \frac{\sigma}{1 - (16d+1)^2 \sigma}.
\]

\begin{theorem}[volume-uniform area law, exact activities;
\texttt{normalized\_exp\_wilson\_loop\_area\_law}]\label{thm:main}
Let $N_c \ge 1$, let $\delta \ge 0$ bound the conjugate-pair activities as
above, and let $t, \varepsilon \ge 0$ and $\sigma \in [0,1]$ satisfy the
explicit window inequalities
\[
  (16d+1)^2\, \rho(\delta)\, e^{t+\varepsilon+1} < 1, \qquad
  16d \cdot
  \frac{\rho(\delta)\, e^{t+\varepsilon+1}}
       {1 - (16d+1)^2 \rho(\delta) e^{t+\varepsilon+1}} \le t,
\]
\[
  (16d+1)^2 \sigma < 1, \qquad
  \rho(\delta)\cdot \exp\!\bigl(16d \cdot K(t;\delta)\bigr) \le \sigma^2 .
\]
Then for every lattice size $N$ and every closed loop $C$,
\[
  \bigl\lVert \langle \tr W_C \rangle_N \bigr\rVert
  \;\le\;
  N_c \cdot
  e^{\loopsupp(C)\cdot 4d \cdot K(t;\delta)} \cdot
  \sigma^{\Area(C)} \cdot
  e^{\loopsupp(C)\cdot 4d \cdot S(\sigma)} .
\]
Every constant on the right is independent of $N$.
\end{theorem}

The hypotheses are pure smallness and geometry inequalities: in an earlier
stage of the development the integrability of the relevant families was
carried as a hypothesis, and it has since been discharged as a theorem
(ledger Addendum~17u; \texttt{normalized\_wilson\_loop\_area\_law\_unconditional}),
so no measure-theoretic side condition remains.  The linearized companion
(activities $1 + f_p$ with $\lVert f_p\rVert$-control) is
\texttt{normalized\_wilson\_loop\_area\_law}; the pipeline is
activity-agnostic and the exact version is obtained from the linearized
proof by the single substitution $2\delta N_c \rightsquigarrow
e^{2\delta N_c}-1$.  Together with the finite-volume pair the programme
comprises four variants:

\begin{center}
\begin{tabular}{lll}
\toprule
Variant & Activities & Lean declaration \\
\midrule
finite volume & linearized & \texttt{finite\_volume\_area\_law} \\
finite volume & exact $\prod_p e^{z_p}$ & \texttt{finite\_volume\_area\_law\_exp} \\
volume-uniform & linearized & \texttt{normalized\_wilson\_loop\_area\_law} \\
volume-uniform & exact $\prod_p e^{z_p}$ & \texttt{normalized\_exp\_wilson\_loop\_area\_law} \\
\bottomrule
\end{tabular}
\end{center}

\begin{theorem}[manifest exponential decay;
\texttt{area\_law\_to\_exp\_area\_decay}]\label{thm:tension}
On any loop family whose perimeter charge is area-subdominant --- i.e.\
$\loopsupp \cdot 4d\,(K + S) \le \lambda \cdot \Area(C)$ for some
$\lambda < -\log\sigma$ --- the bound of Theorem~\ref{thm:main} repackages
as
\[
  \bigl\lVert \langle \tr W_C \rangle_N \bigr\rVert
  \;\le\; N_c \cdot e^{-\tau\,\Area(C)},
  \qquad \tau = (-\log\sigma) - \lambda > 0 ,
\]
a manifest area law with strictly positive string tension $\tau$.
\end{theorem}

\begin{proposition}[machine-checked non-vacuity]\label{prop:nonvacuous}
Two distinct hypothesis windows are involved, and each carries its own
machine-checked witness.
\emph{(i) The repackaging window of Theorem~\ref{thm:tension}.}  Its
quantities $K, S, \sigma, \lambda$ enter only through the inequalities
displayed there --- they are free parameters of the repackaging lemma, not
constrained by the cluster expansion --- and the witness
$\sigma = \tfrac12$, $\lambda = \tfrac12$ yields the strictly positive
tension $\tau = \log 2 - \tfrac12 > 0$
(\texttt{area\_law\_to\_exp\_area\_decay\_window\_nonempty}, proved via
Mathlib's certified numeric bound \texttt{Real.log\_two\_gt\_d9}).  We
emphasize that $\sigma = \tfrac12$ is \emph{not} claimed to lie in the
cluster window of Theorem~\ref{thm:main}, which requires the much smaller
$(16d+1)^2\sigma < 1$.
\emph{(ii) The cluster-expansion smallness window of
Theorem~\ref{thm:main} itself.}  For every $d$ and $N_c$ there exists
$\delta_0 > 0$ such that all activities with $\delta \le \delta_0$ satisfy
the window (lemmas \texttt{clustering\_window\_nonempty} and
\texttt{sun\_clustering\_window\_nonempty}, instantiated at
$t = \varepsilon = 1$; ledger Addendum~17t).
\end{proposition}

\begin{remark}
Making non-vacuity a theorem is a deliberate methodological choice; see
Section~\ref{sec:lean}.
\end{remark}

\section{Proof architecture}\label{sec:architecture}

The proof is a cluster expansion engineered so that every loss is charged
either to the loop's perimeter or to its area, never to the volume.

\paragraph{Loop-tagged expansion and the pinned dichotomy.}
Writing the exact weight as $\prod_p (1 + f_p)$ with
$f_p = e^{z_p} - 1$ and expanding, the numerator becomes a sum over
\emph{pinned} plaquette sets $S_0$ of terms
$\int \tr W_C \cdot \prod_{p \in S_0} f_p \cdot (\text{far factor})$.
Haar selection rules on $SU(N_c)$ kill every pinned term whose plaquette
set cannot fill the loop: the formalized dichotomy
(\texttt{norm\_integral\_exp\_pinned\_term\_le}) bounds each term by
$\mathbb{1}[\Area(C) \le \#S_0] \cdot N_c \cdot \rho(\delta)^{\#S_0}$.
This is where the intrinsic $\mathbb{Z}/N_c$-chain area of
Section~\ref{sec:setting} enters: a surviving monomial exhibits an
algebraic filling of the loop chain inside $S_0$.

\paragraph{Exact activities via shifted Cauchy products.}
For the exact factor the expansion of $\prod_{p\in S_0}(e^{z_p}-1)$ is
organized as a shifted multi-variable Cauchy product
$\sum'_{m} \prod_p z_p^{m_p+1}/(m_p+1)!$ --- every exponent is at least
one, so the occupied plaquette set of each monomial is exactly $S_0$ and
the Haar kill fires uniformly.  The interchange of integral and series is
justified by dominated summability
(\texttt{summable\_prod\_pow\_succ\_div\_factorial}); this is the point at
which the linearized proof generalizes by the substitution
$2\delta N_c \rightsquigarrow e^{2\delta N_c}-1$.

\paragraph{Cancelling the partition function.}
The far factor of a pinned term is a restricted partition function
$Z_{\mathrm{far}}(S_0)$.  A volume-restricted Mayer inversion and a
$Z$-ratio bound
(\texttt{norm\_normalized\_exp\_wilson\_loop\_le\_pinned\_sum}) cancel
$Z_N$ against $Z_{\mathrm{far}}(S_0)$ at a cost exponential in
$\#S_0$ and in the loop's perimeter charge only --- the decisive
volume-uniformity step.

\paragraph{Resumming the pinned gas.}
The remaining sum over pinned sets is controlled by an abstract pinned-gas
estimate (\texttt{sum\_pinned\_dichotomy\_le}) built on the repository's
Koteck\'y--Preiss layer: a 26-module, $\approx$11{,}500-line development
of polymer counting with Penrose-style breadth-first-search spanning-tree
bounds, Ursell coefficients, and the KP convergence
criterion~\cite{Penrose1967,KoteckyPreiss1986}.  The geometric rate
$\sigma$ absorbs one factor per pinned plaquette beyond the area threshold,
yielding the $\sigma^{\Area(C)}$ decay with the perimeter-only prefactors
$e^{\loopsupp\cdot 4d\cdot K}$ and $e^{\loopsupp\cdot 4d\cdot S(\sigma)}$.

\section{The formalization}\label{sec:lean}

\paragraph{Scale and toolchain.}
The area-law dependency cone spans the lattice chain complex (L0), the
Gibbs-measure layer (L1), a 3{,}075-line Wilson-loop monomial module, the
26-module Koteck\'y--Preiss layer, and the 2{,}117-line volume-restricted
gate module containing the headline theorems.  Everything is Lean~4 over
Mathlib and is rebuilt by continuous integration.

\paragraph{Oracle discipline.}
The repository maintains a verification ledger recording, for every
exported theorem, the output of Lean's \texttt{\#print axioms}.  For the
declarations of Sections~\ref{sec:results} the ledger records exactly the
canonical classical trio; e.g.
\begin{lstlisting}
'YangMills.normalized_exp_wilson_loop_area_law'  [propext, Classical.choice, Quot.sound]
'YangMills.normalized_wilson_loop_area_law'      [propext, Classical.choice, Quot.sound]
'YangMills.area_law_to_exp_area_decay'           [propext, Classical.choice, Quot.sound]
\end{lstlisting}
There is no \texttt{sorry} and no project-specific axiom anywhere in this
dependency cone.  For full disclosure: the repository as a whole also
contains clearly labeled exploratory lanes: the experimental modules carry
explicit axioms (fifteen, all confined to \texttt{Experimental/}) and
abandoned proof attempts, while the renormalization-group lane carries
explicitly named \emph{hypotheses} rather than axioms.  None of these is
imported by the area-law chain, and the per-theorem oracle is precisely the
mechanism that makes this separation checkable rather than asserted.

\paragraph{Non-vacuity as a theorem.}
The project's README retains a legacy disclaimer: an early advertised
headline of the shape $\exists\, m > 0$ was discovered to be vacuously
provable and was publicly disclaimed and removed.  The lasting correction
is procedural: every smallness window now comes with a machine-checked
non-emptiness witness (Proposition~\ref{prop:nonvacuous} and ledger
Addenda~17t/20), and integrability side conditions were converted from
hypotheses into theorems (Addendum~17u).  We highlight this as a
transferable practice for formalized analysis: \emph{a conditional theorem
should ship with a formal witness that its conditions are satisfiable}.

\paragraph{Engineering notes for formalizers.}
Constants are carried symbolically ($K$, $S$, $\rho$) rather than
numerically, so that volume-uniformity is syntactically visible in the
statement; the $\int\!\leftrightarrow\!\sum'$ interchanges are localized in
a small number of dominated-convergence lemmas; and the activity-agnostic
design of the pinned pipeline is what reduced the exact-activity upgrade to
a one-parameter substitution --- the entire V4 stage was closed in five
oracle-checked increments (ledger Addenda~18--18d).

\section{Reproducibility and artifacts}\label{sec:repro}

All artifacts are public in the programme repository
\begin{center}
\url{https://github.com/lluiseriksson/THE-ERIKSSON-PROGRAMME}
\end{center}
including the Lean sources (\texttt{YangMills/L1\_GibbsMeasure/RestrictedGate.lean}
for the headline theorems), the verification ledger
(\texttt{docs/VERIFICATION-LEDGER.md}) with per-theorem oracle lines, the
staged plan documents (\texttt{docs/AREA-LAW-VU-PLAN.md}) whose bricks map
one-to-one onto the lemmas of Section~\ref{sec:architecture}, and
continuous-integration builds; the pinned sources and ledger support
independent replay.

\section{Related work and outlook}\label{sec:related}

The strong-coupling area law is classical~\cite{OsterwalderSeiler1978,
Seiler1982}; the contribution here is its machine-checked, explicitly
volume-uniform form for the exact Boltzmann factor.  On the formalization
side, proof assistants have reached deep results in analysis and
probability, but we are not aware of a prior mechanized polymer or cluster
expansion for any lattice field theory.  A companion result in the same
programme --- an exact Catalan evaluation of the factorially weighted
spanning-tree sums that appear in the Koteck\'y--Preiss counting
layer~\cite{Eriksson2026Catalan} --- has been formalized with the same
oracle discipline and is being prepared as generic Mathlib contributions;
sharpening the tree-counting constants of Section~\ref{sec:architecture}
with it is natural future work.  The larger programme --- mass gap,
renormalization-group inputs --- remains open and clearly labeled as such
in the repository; nothing in this paper depends on it.

\subsection*{Acknowledgments}
The Lean formalization was developed interactively with Anthropic's
Claude, in accordance with the ACM authorship policy on AI-assisted work;
every theorem stated here was checked by the Lean~4 kernel, with the axiom
oracle outputs archived in the repository ledger.

\begin{thebibliography}{9}

\bibitem{Wilson1974}
K.~G. Wilson,
\emph{Confinement of quarks},
Phys.\ Rev.\ D \textbf{10} (1974), 2445--2459.

\bibitem{OsterwalderSeiler1978}
K.~Osterwalder and E.~Seiler,
\emph{Gauge field theories on a lattice},
Ann.\ Physics \textbf{110} (1978), 440--471.

\bibitem{Seiler1982}
E.~Seiler,
\emph{Gauge Theories as a Problem of Constructive Quantum Field Theory and
Statistical Mechanics},
Lecture Notes in Physics 159, Springer, 1982.

\bibitem{Penrose1967}
O.~Penrose,
\emph{Convergence of fugacity expansions for classical systems},
in: T.~A.~Bak (ed.), \emph{Statistical Mechanics: Foundations and
Applications}, Benjamin, New York, 1967, 101--109.

\bibitem{KoteckyPreiss1986}
R.~Koteck\'y and D.~Preiss,
\emph{Cluster expansion for abstract polymer models},
Comm.\ Math.\ Phys.\ \textbf{103} (1986), 491--498.

\bibitem{Lean4}
L.~de~Moura and S.~Ullrich,
\emph{The Lean 4 theorem prover and programming language},
in: Automated Deduction -- CADE 28, LNCS 12699, Springer, 2021, 625--635.

\bibitem{mathlib}
The mathlib Community,
\emph{The Lean mathematical library},
in: CPP 2020, ACM, 2020, 367--381.

\bibitem{Eriksson2026Master}
L.~Eriksson,
\emph{The Master Map: An Audit-First, Attack-Resistant Navigation Guide to
the Unconditional Solution of the 4D $SU(N)$ Yang--Mills Existence and Mass
Gap Problem},
ai.viXra:2602.0096, 2026.

\bibitem{Eriksson2026Audit}
L.~Eriksson,
\emph{Mechanical Audit Experiments and Reproducibility Appendix for a
Companion-Paper Programme on 4D $SU(N)$ Yang--Mills Existence and Mass
Gap},
ai.viXra:2602.0117, 2026.

\bibitem{Eriksson2026Catalan}
L.~Eriksson,
\emph{A Machine-Verified Bijective Proof of the Rooted Child-Factorial
Catalan Identity over Spanning Trees of the Complete Graph},
ai.viXra:2607.0001, 2026.

\end{thebibliography}

\end{document}
